\section{Limitations}
\label{sec:limitations}

\paragraph{No ground truth.} The data have no self-reports or think-aloud, so we cannot say which definition is right.
A stall is a stop in net character production, not a judgement that the writer was stuck.

\paragraph{The range varied.} All proportions are of the definitions examined. The stall grid is centred on one
detector's default, no window length has a use outside the detector, and fixed thresholds make up four in ten stall definitions. The
stall detector does not reach pauses shorter than about 15~s, and the pause family covers 2 to 10~s and writer
quantiles; we did not vary thresholds below 2~s, thresholds proportional to within-word speed \cite{guo2020scenario}
or pauses split by linguistic location, because we fixed grids that apply alike to code and essays. We varied the rate per observed hour only; rates per key press or per
word and measures of pause duration, such as mixture components, were not varied, so the recommendation is limited to
rates. We did not compare groups within a dataset.

\paragraph{Differences between datasets.} The datasets differ at once in the task, the recording (CS1 key presses are
reconstructed from edit events) and the length of units, so differences between them are not attributed to a cause.

\paragraph{Intervals and samples.} The primary intervals hold each writer's quantile threshold as detected and leave
out its variation. A resample that redraws it causally left the score and ratio conclusions as they were, but for
$\varphi$ it was shifted above the point values, so how the variation of the threshold bears on $\varphi$ is not
resolved (S2). The default observation intervals depend on the companion study's estimator; every result is also
reported for two segmentations that need only the time of each key press and whether it inserted or deleted a
character.
A definition changes the time in which a stall can be detected (the writer-quantile family
cannot detect before the 100th interval, 3.2--10.3\% of the exposure), which units reach 5 minutes of exposure and, for $\varphi$, which writers
have stalls at all, so definitions differ in their samples as well as their labels. Removing the undetectable time (Section~\ref{sec:detector})
changed no score interval from one side of zero to the other, but moved the CS1 ratio and $\varphi$ in the CS1 2021
assignments above 25\% and left the CS1 2021 blocks uncounted (S2).

\paragraph{Counts and lengths.} The levels of variation in counts could not be separated (S6). The dispersion of
assignments includes differences between assignments, and a $\varphi$ near 1 does not establish a Poisson process: the
window of the detector can make counts underdispersed while heterogeneity makes them overdispersed. The length
decomposition excludes stalls that did not resume, and CS1 days are inferred from gaps of 6~h.

\paragraph{One outcome and its length.} The score association rests on one holistic score in KLiCKe and was not checked
against a second outcome. Its inter-rater agreement ($\kappa$ 0.60 before and 0.76 after discussion) is imperfect, which shrinks
correlations towards zero and can move intervals into the negligible class. The score reflects essay length strongly (0.75), so part of any label's association with it
can run through length; and because the observation time is nearly fixed in KLiCKe, time spent in stalls reduces the
time left for writing, so adjusting for length removes part of what the label counts. Speed and length were adjusted linearly on ranks, and non-linear remainders were not examined.
We did not separate stalls by their position in the session (planning early, review or waiting late under the
fixed 30-minute timer), so the positive corner and the positive length-adjusted values cannot be told apart from that
route. The score association rests on adults recruited online, not on a course. We did not evaluate the labels as
predictors. We say nothing about Japanese input, input method editors, natural work
or other populations.
